\subsection{Transferência de Hohmann Coplanar}
\label{sec:transf_hohmann}

Após a inserção na órbita inicial descrita na Seção~\ref{sec:lancamento_vinculos}, a fase seguinte da missão de RVD/B envolve o ajuste do raio orbital do \emph{chaser}, de modo que ele se aproxime gradualmente da órbita do \emph{target}.
Sendo assim, a transferência de Hohmann permite a passagem de uma órbita circular de raio $r_1$ (do alvo) para outra órbita circular de raio $r_2$ (do alvo). Para isso, assumiu-se que ambas são coplanares e circulares.

Neste trabalho foi considerado que a dinâmica orbital é descrita pelo problema dos dois corpos, sem perturbações externas, as massas são constantes e utilização do parâmetro gravitacional da Terra ($\mu$).
A descrição do problema dos dois corpos é realizada no referencial ECI. A equação fundamental se dá a partir da equação do movimento orbital absoluto:

\begin{equation}
\ddot{\vec r} = -\frac{\mu}{r^3}\vec r,
\qquad r = \|\vec r\|,
\end{equation}

obtém-se a expressão da energia orbital específica:

\begin{equation}
\varepsilon = \frac{v^2}{2} - \frac{\mu}{r},
\end{equation}

ainda, para uma órbita kepleriana de semieixo maior $a$, é assumido o valor:

\begin{equation}
\varepsilon = -\frac{\mu}{2a}.
\end{equation}

Ao igualar as duas formas da energia e isolar $v^2$, obtém-se à equação vis-viva:

\begin{equation}
v^2 = \mu \left(\frac{2}{r} - \frac{1}{a}\right).
\end{equation}

%%%%%%%%%%%%%%%%%%%%%%%%%%
%%% ETAPAS DA MANOBRA
%%%%%%%%%%%%%%%%%%%%%%%%%%
\subsubsection{Etapas da manobra}

A manobra de Hohmann é composta por dois impulsos.
O primeiro impulso ($\Delta v_1$) é aplicado no perigeu da órbita inicial, levando o veículo de $r_1$ para uma órbita elíptica de transferência com semieixo maior:

\begin{equation}
a_T = \frac{r_1+r_2}{2}.
\end{equation}

A variação de velocidade no perigeu é dado por:

\begin{equation}
\Delta v_1 = \sqrt{2\mu \left(\frac{1}{r_1} - \frac{1}{2a_T}\right)} - \sqrt{\frac{\mu}{r_1}}.
\end{equation}

O segundo impulso ($\Delta v_2$) é aplicado no apogeu da órbita de transferência para circularizar a trajetória em $r_2$:

\begin{equation}
\Delta v_2 = \sqrt{\frac{\mu}{r_2}} - \sqrt{2\mu \left(\frac{1}{r_2} - \frac{1}{2a_T}\right)}.
\end{equation}

O tempo de voo na órbita de transferência é metade do período da órbita elíptica:

\begin{equation}
t_T = \pi \sqrt{\frac{a_T^3}{\mu}}.
\end{equation}

Por fim, a variação total de velocidade é a soma dos dois impulsos:

\begin{equation}
\Delta v_{H} = \Delta v_1 + \Delta v_2.
\end{equation}

\subsubsection{Condição de fase e janela de lançamento}

Para que a manobra seja considerada viável, o alvo deve estar à frente do veículo visitante, de acordo com o referencial ECI, por um ângulo inicial compatível com o tempo de transferência $t_{Transf}$.
% Esse ângulo é obtido a partir da diferença entre as velocidades angulares médias das duas órbitas.

\begin{equation}
\varphi =
\pi - \left(\bar{n}_2 - \bar{n}_1\right)\,t_{Transf},
\end{equation}